\documentclass[a4paper]{article}
% Español (México) con XeLaTeX: fontspec para fuentes Unicode, polyglossia
% para la división silábica y los nombres del documento en la variante mexicana.
\usepackage{fontspec}
\usepackage{polyglossia}
\setdefaultlanguage[variant=mexican]{spanish}
% Los nombres chinos van como «pinyin (original)»: xeCJK da a los caracteres
% Han su propia fuente; sin ella XeLaTeX los omite con un aviso.
% FandolSong no cubre algunos caracteres de nombres (U+73FA); WenQuanYi Zen Hei sí.
\usepackage[AutoFallBack=true]{xeCJK}
\setCJKmainfont{FandolSong-Regular.otf}
\setCJKfallbackfamilyfont{\CJKrmdefault}{WenQuanYi Zen Hei}
% polyglossia no traduce los nombres de listings.
\AtBeginDocument{\providecommand{\lstlistingname}{}\renewcommand{\lstlistingname}{Listado}%
  \providecommand{\lstlistlistingname}{}\renewcommand{\lstlistlistingname}{Índice de listados}}
\usepackage{amsmath, amssymb}
\usepackage{graphicx}
\usepackage[margin=2.5cm]{geometry}
\usepackage[most]{tcolorbox}
\usepackage{etoolbox}
\usepackage{listings}
\usepackage{hyperref}
\usepackage{booktabs}
\usepackage{float}
\newtcolorbox{knowledgebox}[1]{enhanced, colback=blue!5!white, colframe=blue!75!black, colbacktitle=blue!75!black, coltitle=white, fonttitle=\bfseries, title={#1}, attach boxed title to top left={yshift=-2mm, xshift=2mm}, boxrule=1pt, sharp corners}
\newtcolorbox{importantbox}[1]{enhanced, colback=yellow!10!white, colframe=yellow!80!black, colbacktitle=yellow!80!black, coltitle=black, fonttitle=\bfseries, title={#1}, sharp corners}
\newtcolorbox{warningbox}[1]{enhanced, colback=red!5!white, colframe=red!75!black, colbacktitle=red!75!black, coltitle=white, fonttitle=\bfseries, title={#1}, sharp corners}
\lstset{language=Python, basicstyle=\ttfamily\small, keywordstyle=\color{blue}, stringstyle=\color{red!60!black}, commentstyle=\color{green!60!black}, breaklines=true, frame=single, numbers=left, numberstyle=\tiny\color{gray}, captionpos=b, extendedchars=true}
\newcommand{\notetitle}{Reward Model y modelado probabilístico y estadístico}
\newcommand{\noteauthors}{Elaboradas a partir de materiales públicos del curso}
\newcommand{\notedate}{2025}
\newcommand{\videochannel}{Wudaokou Nashi (五道口纳什)}
\newcommand{\videopublishdate}{2025}
\newcommand{\videoduration}{aproximadamente 20 minutos}
\newcommand{\videourl}{}
\newcommand{\videocoverpath}{cover.jpg}
\newcommand{\repourl}{https://github.com/hqhq1025/ai-course-notes}

% Footer with repo info
\usepackage{fancyhdr}
\pagestyle{fancy}
\fancyhf{}
\fancyfoot[C]{\thepage}
\fancyfoot[R]{\footnotesize\color{gray}\href{\repourl}{AI Course Notes}}
\renewcommand{\headrulewidth}{0pt}
\renewcommand{\footrulewidth}{0pt}

\begin{document}
\begin{titlepage}\centering
{\Large Notas del curso\par}\vspace{1.2cm}{\huge\bfseries \notetitle\par}\vspace{0.8cm}{\large \noteauthors\par}\vspace{0.3cm}{\large \notedate\par}\vspace{1.2cm}
\ifdefempty{\videocoverpath}{}{\includegraphics[width=0.82\textwidth,height=0.45\textheight,keepaspectratio]{\videocoverpath}\par}
\vfill
\begin{tcolorbox}[width=0.9\textwidth, colback=black!2!white, colframe=black!60, sharp corners]
\textbf{Autor o canal del video}: \videochannel\par\textbf{Fecha de publicación}: \videopublishdate\par\textbf{Duración del video}: \videoduration\par
\ifdefempty{\videourl}{}{\textbf{Enlace del video}: \href{\videourl}{\nolinkurl{\videourl}}\par}\par
\textbf{Página del proyecto}: \href{\repourl}{\nolinkurl{\repourl}}\par
\end{tcolorbox}
\end{titlepage}
\tableofcontents\newpage

\section{Introducción}
En esta sesión se retoma el Reward Model mediante el modelo de Bradley-Terry, para explicar la relación entre el modelado probabilístico y estadístico, el aprendizaje profundo y la estimación de máxima verosimilitud.

\begin{importantbox}{Relación entre los tres}
\begin{enumerate}
  \item \textbf{Modelado probabilístico y estadístico}: define el proceso de generación de los datos (por ejemplo, el modelo de preferencias de Bradley-Terry)
  \item \textbf{Estimación de máxima verosimilitud (MLE)}: define la función de pérdida a partir del modelo estadístico
  \item \textbf{Aprendizaje profundo}: sirve como herramienta de resolución, optimizando la pérdida mediante descenso de gradiente
\end{enumerate}
\end{importantbox}

\section{Modelo de preferencia de Bradley-Terry}
Dadas dos response $y_w$ (preferred) y $y_l$ (rejected), el modelo de Bradley-Terry define la probabilidad de preferencia:
$$P(y_w \succ y_l) = \sigma(r(y_w) - r(y_l))$$
donde $r(\cdot)$ es la salida del Reward Model y $\sigma$ es la función sigmoid.

\section{De MLE a Loss}
Maximizar la verosimilitud $\Leftrightarrow$ minimizar la log-verosimilitud negativa:
$$\mathcal{L} = -\log \sigma(r(y_w) - r(y_l))$$

\subsection{Resumen de la sección}
El entrenamiento del Reward Model es, en esencia, modelado probabilístico-estadístico más resolución por MLE. El aprendizaje profundo es solo la herramienta de optimización.

\newpage
\section{Aviso de modelado: no te fijes solo en la envoltura de la red neuronal}
Esta clase insiste una y otra vez en Bradley-Terry y en MLE; en realidad, le recuerda al investigador que el Reward Model es, ante todo, un modelo probabilístico y, solo después, una tarea de entrenamiento de una red neuronal. Si la atención se centra únicamente en la arquitectura de la red y en el ajuste del loss, es fácil pasar por alto si se cumplen los supuestos estadísticos de los datos de preferencia.

\begin{warningbox}{Errores comunes del Reward Model}
\begin{itemize}
  \item Tratar el modelado de preferencias completamente como un problema de clasificación de caja negra
  \item No verificar si los datos realmente satisfacen el supuesto de preferencia por pares
  \item Fijarse solo en el loss de entrenamiento, sin analizar si la distribución del reward es estable
\end{itemize}
\end{warningbox}

\subsection{Resumen de la sección}
Para dominar el Reward Model, lo esencial es primero entender cómo se modela la probabilidad de preferencia, y después decidir qué herramienta de aprendizaje profundo se usará para ajustarla.

\newpage
\section{Síntesis y ampliación}
\begin{enumerate}
  \item Dominar el modelado probabilístico y estadístico es la base para hacer bien investigación en RL
  \item Bradley-Terry es el método de modelado estándar para el Reward Model
  \item El aprendizaje profundo es la herramienta para resolver, el modelo probabilístico es el núcleo
\end{enumerate}
\end{document}
